\documentclass[11pt]{article}
\usepackage[margin=1in]{geometry}
\usepackage{amsmath,amssymb,amsthm,hyperref}
\newtheorem{theorem}{Theorem}[section]
\newtheorem{lemma}[theorem]{Lemma}
\newtheorem{corollary}[theorem]{Corollary}
\newcommand{\R}{\mathbb R}
\newcommand{\tr}{\operatorname{tr}}
\newcommand{\rank}{\operatorname{rank}}
\newcommand{\norm}[1]{\left\lVert#1\right\rVert}
\renewcommand{\labelitemi}{$\bullet$}
\title{A charged relative fractional graph update using feasible flows}
\author{Independent dynamic research session}
\date{November 15, 2026}
\begin{document}
\maketitle

\paragraph{Status and scope.}
This is a proposed constructive theorem with a complete exact-real
arithmetic argument, including factor formation and a small relative
Gram computation. It uses the feasible electrical-flow guarantee of
Kelner--Orecchia--Sidford--Zhu, not merely the usual shifted-potential
energy guarantee. Independent checks within this session verified the
flow, clipping, and Gram-compression steps. Historical priority and
finite-word stability have not been established.

For general fixed $\alpha$, scalar evaluations of positive real powers
are charged as scalar operations, as in the source graph-update model.
The quadrature coefficients are computed to the same error budgets.
For $\alpha=1/2$, one can instead take $h=\log(2)/H$ with dyadic integer
$H=O_\alpha(\mathcal L)$, use nodes $2^{j/H}$, and construct them by
repeated square roots; scalar constant approximations use convergent
series. The small-matrix Jacobi step below also uses only arithmetic,
comparisons, and positive square roots. Thus no exact general matrix
spectral-decomposition oracle is part of the arithmetic model.

\section{Statement}
Let $L$ be the combinatorial Laplacian of a connected weighted undirected
graph with $n$ vertices and $m$ edges. Vertex masses are one. Let
\[
 L'=L+d b b^T,\qquad d>0,\qquad b=e_p-e_q,\qquad
 \tau=d b^TL^\dagger b>0,
\]
where insertion and increasing an existing edge are both allowed.
Fix $0<\alpha<1$, and put $B_0=\alpha\log(1+\tau)$.
Every inverse power below is zero on constants. Set
\[
 \Delta=L'^\alpha-L^\alpha,\qquad
 C=L^{-\alpha/2}\Delta L^{-\alpha/2}.
\]
The known determinant identity is
$C\succeq0$, $\log\det_{\mathbf1^\perp}(I+C)=B_0$.

\begin{theorem}[Explicit relative factor, with charged construction]
\label{thm:main}
For $0<\varepsilon,\delta<1/2$, there is a randomized exact-real
arithmetic algorithm constructing a matrix $X\in\R^{n\times r}$ with
$X^T\mathbf1=0$ such that, with probability at least $1-\delta$,
\begin{equation}\label{eq:answer}
 -\frac\varepsilon8 L^\alpha
 \preceq \Delta-XX^T
 \preceq\frac{3\varepsilon}4 L^\alpha.
\end{equation}
In particular its two-sided relative operator error is at most
$\varepsilon$. On the same success event, its rank obeys
\begin{equation}\label{eq:rank}
 r<\frac{B_0}{\log(1+\varepsilon/4)}
 \quad\hbox{if }r>0.
\end{equation}
Let $\mathcal L=1+\log\bigl(2+n(1+\tau)/\varepsilon\bigr)$.
There are raw width $q=O_\alpha(\mathcal L^2)$ and Gram quadrature
width $N=O_\alpha(\mathcal L^2)$ such that construction uses
\begin{itemize}
\item expected $O(q+1)$ feasible electrical-flow calls on graphs with
$n+1$ vertices and $m+n$ edges;
\item expected $O(qN+q+1)$ ordinary SDD potential calls;
\item $O(Nnq^2+q^4\mathcal L+nqr)$ further arithmetic operations, including
all raw and final factor entries.
\end{itemize}
These counts include the optional certified leverage estimation and
precede confidence amplification. With a supplied leverage bound the
raw construction has exactly $q$ flow targets; amplification adds the
logarithmic repetitions described below.
All required solve and flow tolerances have
$\log(1/\eta)=O_\alpha(\mathcal L)$, before logarithmic confidence
amplification. The resulting expected work is
\[
 \widetilde O_\alpha(m+n),
\]
where suppressed factors depend only logarithmically on
$n,1+\tau,1/\varepsilon,1/\delta$. There is no suppressed dependence
on $\lambda_2(L)$, $\norm L$, conductance ratios, or a graph condition
number. A constant-factor estimate of $\tau$ suffices and is computable
within the same work bound.
\end{theorem}

For any fixed $0<\rho<1$, the same construction, with different error
allocations, returns, on the same success event, relative error at most
$\varepsilon$ and
$r<B_0/\log(1+(1-\rho)\varepsilon)$ when $r>0$. Thus the leading
rank constant can approach the sharp determinant count. Dependence on
$1/\rho$ is only logarithmic in construction work.

The theorem is a two-sided relative approximation by a PSD factor.
It does not assert $XX^T\preceq\Delta$. Per edit it may reconstruct
solver data from the current graph; no sublinear dynamic-update bound
or uncharged initial dense fractional matrix is being asserted.

\section{Graph-specific scale control}
Scale all conductances by $1/d$. It suffices to work with $d=1$,
$u=b$, $\norm u^2=2$, and return $d^{\alpha/2}X$ at the end.
Write $A=L/d$ in this normalization and $v=A^\dagger u$.
The graph maximum principle puts every coordinate of $v$ between
$v_q$ and $v_p$, and $v_p-v_q=\tau$. Thus
\begin{equation}\label{eq:moments}
 \norm v^2\le n\tau^2,\qquad
 u^TA^{-\alpha}u\le\tau^\alpha 2^{1-\alpha}.
\end{equation}
The second inequality is spectral H\"older interpolation.
These estimates are special to a dipole on a unit-mass graph.

Let $R_t=(A+tI)^{-1}$ on the whole vertex space and
$s(t)=u^TR_tu$. The exact relative integral is
\begin{equation}\label{eq:integral}
 C=c_\alpha\int_0^\infty
 \frac{t^\alpha A^{-\alpha/2}R_tuu^TR_tA^{-\alpha/2}}
 {1+s(t)}\,dt,\qquad c_\alpha=\frac{\sin(\pi\alpha)}\pi.
\end{equation}
In particular $0\le s(t)\le\tau$.

\begin{lemma}[Relative integral tails without a spectral endpoint]
For $T>0$, the traces of the portions of \eqref{eq:integral} below
and above $T$ are bounded respectively by
\[
 K_\alpha T\norm v^2,\qquad
 K_\alpha T^{\alpha-1}\tau^\alpha2^{1-\alpha}.
\]
\end{lemma}
\begin{proof}
Drop the denominator. For the low tail use
\[
 \sup_{\lambda>0}
 \frac{\lambda^{2-\alpha}}{(\lambda+t)^2}
 =K'_\alpha t^{-\alpha},
\]
so the trace integrand is at most $K_\alpha\norm{A^\dagger u}^2$.
For the high tail use $R_t^2\preceq t^{-2}I$ and
\eqref{eq:moments}, then integrate $t^{\alpha-2}$.
\end{proof}

Choose a sufficiently large constant $K_\alpha$ and cutoffs
\begin{equation}\label{eq:cutoff}
 t_- = \frac{\varepsilon}{K_\alpha n(1+\tau)^2},\qquad
 t_+ = \left(\frac{K_\alpha(1+\tau)}\varepsilon\right)^{1/(1-\alpha)}.
\end{equation}
Their logarithms have magnitude $O_\alpha(\mathcal L)$, and each
omitted relative tail has norm at most $\varepsilon/100$ after enlarging
the constant. No estimate of an extreme eigenvalue is used.

\begin{lemma}[Positive logarithmic quadrature]
There is a logarithmic trapezoid rule with $q=O_\alpha(\mathcal L^2)$,
nodes in $[\exp(-O_\alpha(\mathcal L)),
\exp(O_\alpha(\mathcal L))]$, and positive weights
\[
 a_j=\frac{c_\alpha h t_j^{\alpha+1}}{1+s(t_j)}
\]
whose raw exact factor $Z_0=[\sqrt{a_j}R_{t_j}u]_j$ obeys
\[
 \norm{A^{-\alpha/2}Z_0Z_0^TA^{-\alpha/2}-C}
 \le\varepsilon/32.
\]
Moreover $J:=qh=O_\alpha(\mathcal L)$ and
$\norm{A^{-\alpha/2}Z_0}_F^2\le K_\alpha\tau J$.
\end{lemma}
\begin{proof}
Put $t=e^y$. The transformed matrix integrand is analytic in any fixed
strip $|\operatorname{Im}y|\le a<\pi/2$. Here
$\operatorname{Re}s(e^y)\ge0$, so $|1+s(e^y)|\ge1$.
Also $|\lambda+te^{i\theta}|\ge\cos(\theta/2)(\lambda+t)$.
The integral of the norm on either strip boundary is at most a constant
times
\[
 c_\alpha\int_0^\infty t^\alpha
       u^T R_t A^{-\alpha}R_tu\,dt=\alpha\tau.
\]
The scalar equality follows eigenvalue by eigenvalue. Shifting the
Fourier contour and summing the sampled nonzero frequencies bounds
infinite-trapezoid error by
$K_\alpha\tau/(e^{2\pi a/h}-1)$.
Choose $h^{-1}=O_\alpha(\mathcal L)$ sufficiently large.
The same low and high envelopes from the preceding lemma bound the
discarded grid tails by geometric sums. Rounding the cutoffs by at
most a constant number of grid points and enlarging constants proves
the error and width statements.
Finally every transformed positive integrand has trace at most
$K_\alpha\tau$, since
$t^{\alpha+1}\lambda^{1-\alpha}/(\lambda+t)^2\le K_\alpha$.
Summing the $q$ node traces proves the Frobenius bound.
\end{proof}

\section{Constructing relative-accurate resolvent columns}
An arbitrary shifted-potential energy error can contaminate an
arbitrarily small positive eigenmode of $A$. The following construction
uses feasible flows to avoid that failure.

\begin{lemma}[Clipping a dipole potential]
Suppose $\widetilde v$ has
$\norm{\widetilde v-v}_A\le\eta\sqrt\tau$, $0<\eta<1$.
Clip every coordinate of $\widetilde v$ into the interval with endpoints
$\widetilde v_q,\widetilde v_p$, then subtract its mean to obtain
$\widehat v$. Then
\begin{equation}\label{eq:clip}
 \norm{\widehat v-v}_A\le\eta\sqrt\tau,\qquad
 \norm{\widehat v}\le\sqrt n(1+\eta)\tau.
\end{equation}
\end{lemma}
\begin{proof}
Cauchy--Schwarz in the $A$ energy gives
$|u^T(\widetilde v-v)|\le\eta\tau$, so the endpoint order is correct
and the interval length is at most $(1+\eta)\tau$.
Common scalar clipping is nonexpansive on every edge difference and
fixes both endpoint values. It decreases $x^TAx$ and preserves $u^Tx$.
Therefore it cannot increase
$\norm{x-v}_A^2=x^TAx-2u^Tx+\tau$.
Subtracting the mean changes neither expression. The coordinate range
bound gives the Euclidean bound.
\end{proof}

Orient the graph, write $A=D^TWD$, and add a ground vertex with an
edge of conductance $t$ from each original vertex. For demand
$(\widehat v,0)$ a feasible flow consists of original-edge flows $f$
and ground flows $g$ with
\[
 D^Tf+g=\widehat v,
 \qquad \mathcal E(f,g)=f^TW^{-1}f+t^{-1}\norm g^2.
\]
The minimizing flow has $y=R_t\widehat v$,
$f_*=WDy$, $g_*=ty$, and energy
$E_*=\widehat v^TR_t\widehat v$.

\begin{lemma}[Feasible-flow resolvent column]\label{lem:flow}
Given a feasible flow of energy at most $(1+\zeta^2)E_*$, output
$\widehat x_t=D^Tf=\widehat v-g$. Then
\begin{align}
 \norm{A^{-\alpha/2}(\widehat x_t-R_tu)}
 &\le K_\alpha\eta\sqrt\tau\,t^{-(1+\alpha)/2}
       +\zeta\norm{\widehat v}\,t^{-\alpha/2}.\label{eq:flowerr}
\end{align}
Also, for $\zeta\le1$,
\begin{equation}\label{eq:flowmom}
 f^TW^{-1}f\le2\norm{\widehat v}^2/t,
 \qquad \norm{\widehat x_t}^2\le6\norm{\widehat v}^2.
\end{equation}
\end{lemma}
\begin{proof}
Feasible-flow optimality gives exact orthogonality, so the energy excess
is
$E=\norm{f-f_*}_{W^{-1}}^2+t^{-1}\norm{g-g_*}^2\le\zeta^2E_*$.
For $e=D^T(f-f_*)=-(g-g_*)$, Thomson's variational principle and
feasibility give $\norm e_{A^\dagger}^2\le E$ and
$\norm e^2\le tE$. Spectral H\"older interpolation implies
\[
 \norm{A^{-\alpha/2}e}^2
 \le(\norm e_{A^\dagger}^2)^\alpha(\norm e^2)^{1-\alpha}
 \le t^{1-\alpha}E
 \le\zeta^2t^{-\alpha}\norm{\widehat v}^2.
\]
The exact original-edge divergence is $A R_t\widehat v$ and
$A R_t v=R_tu$. The remaining bias is bounded using
\[
 \norm{A^{-\alpha/2}A R_t(\widehat v-v)}^2
 \le\sup_{\lambda>0}
 \frac{\lambda^{1-\alpha}}{(\lambda+t)^2}
 \norm{\widehat v-v}_A^2
 \le K_\alpha t^{-1-\alpha}\eta^2\tau.
\]
For the moments, total energy is at most
$2E_*\le2\norm{\widehat v}^2/t$.
Also $\norm g^2\le2\norm{\widehat v}^2$ and
$\widehat x_t=\widehat v-g$, giving \eqref{eq:flowmom}.
\end{proof}

\paragraph{The actual graph primitives.}
A standard SDD solver constructs $\widetilde v$ with the energy
guarantee above. The feasible flow is supplied by the electrical-flow
algorithm of \cite[Definition 2.4, Theorems 3.2 and 4.1]{KOSZ}.
Its expected excess-energy guarantee at parameter $\zeta^2/4$ gives
the desired bound with constant probability by Markov's inequality.
Independent repetitions and selection of the least-energy feasible
flow give arbitrary confidence: all candidates have the same unknown
optimum, so least energy also means least excess. The cycle algorithm
preserves feasibility exactly. Its stored non-tree flows and residual
demands recover tree flows by subtree sums in $O(m+n)$ work. The
arithmetical runtime is nearly linear with logarithmic accuracy cost,
independent of conductance ratios.

Compute each scalar $s(t_j)$ separately by an ordinary shifted SDD
solve at energy accuracy $\kappa$. Then
$\widehat s_j=\max(0,u^T\widetilde x_j)$ satisfies
$|\widehat s_j-s(t_j)|\le\kappa s(t_j)$.
For $\kappa\le1/2$, replacing the denominator by $1+\widehat s_j$
changes its square-root weight by relative $O(\kappa)$.
Construct
\[
 Z_j=\sqrt{\frac{c_\alpha h t_j^{\alpha+1}}{1+\widehat s_j}}
         \,\widehat x_{t_j}.
\]
All these columns are mean zero, by flow feasibility and zero ground
demand. From \eqref{eq:flowerr}, \eqref{eq:clip}, and the quadrature
Frobenius bound, sufficiently small tolerances with
\[
 \log(1/\eta),\log(1/\zeta),\log(1/\kappa)
       =O_\alpha(\mathcal L)
\]
make
\begin{equation}\label{eq:raw}
 \norm{A^{-\alpha/2}ZZ^TA^{-\alpha/2}-C}\le\varepsilon/8.
\end{equation}
For clarity, a sufficient way to allocate them is to ensure
\[
 K_\alpha\sqrt{\tau J}\,(\eta+\kappa)
 +K_\alpha\zeta\sqrt{n\tau^2Jt_+}
 \le \frac{\varepsilon}{K_\alpha(1+\sqrt{\tau J})}.
\]
The inequality $\norm{TT^T-SS^T}\le
(2\norm S+\norm{T-S})\norm{T-S}$ then bounds factor error;
constants can allocate $\varepsilon/32$ to quadrature and the rest to
column and scalar errors. Every parameter on this display has
logarithm $O_\alpha(\mathcal L)$.

The returned flow energies also provide a computable certificate
$H_0\ge\tr Z^TA^\dagger Z$, and $N_0=\norm Z_F^2$ is directly
computable. Specifically, sum $f_j^TW^{-1}f_j$ times their squared
column weights. The moment bounds yield
\begin{equation}\label{eq:rawmom}
 H_0\le K_\alpha n\tau^2t_+^\alpha,
 \qquad N_0\le K_\alpha n\tau^2t_+^{\alpha+1}.
\end{equation}
Factors $h\sum t_j^a\le K_a t_+^a$ were absorbed in constants.

\section{Small-Gram compression with charged scalar solves}
The raw width is logarithmic in $n$, but may exceed the determinant
budget rank. It can be compressed without a fractional vector action.
Compute the $q$-by-$q$ matrix
\begin{equation}\label{eq:gram}
 H=Z^TA^{-\alpha}Z
   =c_\alpha\int_0^\infty s^{-\alpha}Z^T(A+sI)^{-1}Z\,ds.
\end{equation}
The low and high tails have norm at most
\[
 K_\alpha H_0 S_-^{1-\alpha},\qquad
 K_\alpha N_0 S_+^{-\alpha},
\]
respectively. In normalized units choose, for a sufficiently large
constant, for example
\[
 S_- =\left(\frac{\varepsilon}{K_\alpha(1+H_0)}\right)^{1/(1-\alpha)},
 \qquad
 S_+ =\left(\frac{K_\alpha(1+N_0)}\varepsilon\right)^{1/\alpha}.
\]
By \eqref{eq:rawmom}, both logarithms have magnitude
$O_\alpha(\mathcal L)$.
The logarithmic integrand has a fixed analytic strip and trace envelope
\[
 K_\alpha\min(H_0s^{1-\alpha},N_0s^{-\alpha}).
\]
Its integrated envelope is at most
$K_\alpha H_0^\alpha N_0^{1-\alpha}$, a polynomial in
$n,1+\tau,1/\varepsilon$. The same trapezoid argument therefore
gives $N=O_\alpha(\mathcal L^2)$ nodes for additive error
$\varepsilon/32$.

At every Gram node solve $(A+sI)y_j=Z_j$ by the ordinary SDD energy
solver, and form and symmetrize the matrix $Z^T\widehat Y$.
If every column has energy relative error at most $\theta$, its Gram
error is at most
\[
 \theta\tr Z^T(A+sI)^{-1}Z.
\]
Indeed its $(i,j)$ entry is bounded by
$\theta\sqrt{Z_i^TR_sZ_i}\sqrt{Z_j^TR_sZ_j}$; sum squares.
The integrated envelope bounds accumulated solve error using
$\log(1/\theta)=O_\alpha(\mathcal L)$.
Compute the quadrature and solves accurately enough that their total
error is at most $\varepsilon/16$. Approximately diagonalize the
resulting symmetric matrix as described below to a residual at most
$\varepsilon/16$, and call the resulting explicitly orthogonally
diagonalized matrix $\widehat H$. Then
\begin{equation}\label{eq:gramerr}
 \norm{\widehat H-H}\le\varepsilon/8.
\end{equation}
Each node requires $q$ solves and $O(nq^2)$ arithmetic to form its
Gram. No factor whitening or inverse fractional-power action is used.

Take an orthonormal matrix $U$ spanning the eigenvectors of
$\widehat H$ with eigenvalues greater than $\varepsilon/2$, and
return $X=ZU$. This does not require an exact algebraic eigensolver.
Repeatedly choose the largest off-diagonal entry of the current small
symmetric matrix and apply the exact two-dimensional Jacobi rotation,
using arithmetic and positive square roots. Each rotation reduces the
squared off-diagonal Frobenius norm by at least the fraction
$2/[q(q-1)]$. Stop when the off-diagonal Frobenius norm is at most
$\varepsilon/16$ (equivalently, its square is at most
$\varepsilon^2/256$),
and keep the accumulated orthogonal matrix and the remaining diagonal.
Discarding the off-diagonal residual defines $\widehat H$ explicitly.
Scanning for the largest entry costs $O(q^2)$ per rotation, so
$O(q^4\log(2+\norm{H}_F/\varepsilon))=O_\alpha(q^4\mathcal L)$
arithmetic operations suffice, with no eigenvalue-separation assumption.
For $q=1$ this is immediate. Then
\[
 0\preceq ZZ^T-XX^T,\qquad
 \norm{A^{-\alpha/2}(ZZ^T-XX^T)A^{-\alpha/2}}
 \le\varepsilon/2+\varepsilon/8=5\varepsilon/8.
\]
Combining with \eqref{eq:raw} gives \eqref{eq:answer}.
By Weyl's inequality, every retained eigenvalue forces an eigenvalue
of $H$ above $3\varepsilon/8$, hence one of $C$ above
$\varepsilon/4$. The nonzero eigenvalues of $H$ and
$A^{-\alpha/2}ZZ^TA^{-\alpha/2}$ coincide. The determinant budget
therefore proves \eqref{eq:rank}.

All primitive success events are amplified and union bounded.
Ordinary potential solves can use energy-distance median selection;
feasible flows use least-energy selection. Their confidence overhead
is logarithmic in $(qN+q+1)/\delta$. Forming $ZU$ costs $O(nqr)$.
The displayed primitive counts and algebra costs prove the construction
bound. All normalization, sparse multiplication, flow extraction, and
factor entries are included.

To make the work bound unconditional on successful solves, schedule the
Gram cutoffs and accuracies from the a priori upper bounds on the
right side of \eqref{eq:rawmom}, rather than from arbitrarily large
unsuccessful solver outputs. If a computed moment certificate exceeds
those bounds, declare failure and stop. Likewise cap the number of
Jacobi rotations by the bound implied by these polynomial moments;
failure to reach its residual target declares failure. Neither cap
fires on the simultaneous success event used in the proof. These
guards preserve the stated expected work even on the exceptional
outputs of graph primitives. A promised constant-factor bound on the
fixed $\tau$ schedules all budgets. It can also be obtained by a
certified Las Vegas preliminary procedure: compute a feasible flow
$f$ for demand $u$ and an approximate potential $w$, and evaluate
\[
 D_0=2u^Tw-w^TAw\le\tau\le E_0=f^TW^{-1}f.
\]
Repeat independent constant-accuracy calls until
$0<D_0$ and $E_0\le2D_0$. Standard constant-success potential and
flow bounds make this event have probability bounded below by a
positive constant. Every accepted interval is valid by weak duality,
so $E_0$ is a certified upper bound between $\tau$ and $2\tau$,
with expected nearly-linear preliminary work. Consequently failed
resistance estimates cannot cause unbounded subsequent accuracy logs.

To obtain the near-sharp variant, take raw error and total Gram error
each at most $\rho\varepsilon/4$ and retain diagonalized Gram values
above $(1-\rho/2)\varepsilon$. The final relative error is at most
$\varepsilon$, while every retained value forces a true correction
eigenvalue above $(1-\rho)\varepsilon$. All accuracy logarithms merely
gain $O(\log(1/\rho))$.

\section{Preconditioning and limits}
Let $M=L^\alpha+XX^T$. Since $M\succeq L^\alpha$, the two-sided
guarantee implies
\[
 (1-\varepsilon)M\preceq L'^\alpha\preceq(1+\varepsilon)M.
\]
Thus $M$ is a relative spectral preconditioner. This statement concerns
its representation and quality. Applying $M^{-1}$ by Woodbury still
requires a base $L^{-\alpha}$ action. If that action costs $T_{\rm base}$,
the additional core setup is
$rT_{\rm base}+O(nr^2+r^3)$ and each application costs
$T_{\rm base}+O(nr+r^2)$. No base inverse action is free here.

For fixed leverage the factor construction itself is nearly linear
up to logarithms and has optimal-order worst-case rank $B_0/\varepsilon$.
The logarithmic raw-width bound is consistent with the sharp relative
lower examples: realizing many separated relative modes with a single
unit-mass edge can require extremely large graph dimension. There is
no arbitrary-SPD analogue of the dipole maximum-principle estimate.
The restricted noisy-resolvent barrier in the companion session note
therefore does not contradict this explicit-graph construction.

\begin{thebibliography}{9}
\bibitem{KOSZ} J. A. Kelner, L. Orecchia, A. Sidford, Z. Allen Zhu,
\emph{A Simple, Combinatorial Algorithm for Solving SDD Systems in
Nearly-Linear Time}, STOC 2013.
\url{https://arxiv.org/abs/1301.6628}.
Primary complete paper:
\url{https://math.mit.edu/~kelner/publications/ASimpleCombinatorialAlgorithmforSolvingSDDSystemsinNearly-LinearTime.pdf}.
\bibitem{KMP} I. Koutis, G. L. Miller, R. Peng,
\emph{A nearly-$m\log n$ time solver for SDD linear systems}, 2011.
\url{https://arxiv.org/abs/1102.4842}.
\bibitem{BCKS} B. Beckermann, A. Cortinovis, D. Kressner, M. Schweitzer,
\emph{Low-rank updates of matrix functions II: Rational Krylov methods},
SIAM J. Numer. Anal. 59 (2021), 1325--1347.
\url{https://arxiv.org/abs/2008.11501}.
\bibitem{Repo} Repository manuscripts
\texttt{spectral-gap/relative-fractional-response.tex} and
\texttt{matmul/graph-fractional-updates.tex}, read November 15, 2026.
\end{thebibliography}
\end{document}
